\chapter{Fase 4: el Radar, ¿dónde hay espacio?}
\label{cap:d-f4}
\textbf{Fechas}: 28 y 29 de septiembre de 2026. \textbf{Notas}: \codigo{08-radar.md}, \codigo{09-auditoria-fases-1-4.md}.

\textbf{Qué se buscaba}: un índice de presión de 0 a 1 por lugar y mes, un estado en palabras (tranquilo, concurrido,
saturado) y una predicción del estado futuro.

\begin{opciones}[4.1 · Los pesos del índice · 28-sep-2026]
\textbf{Evidencia que motivó la pregunta}: el sur solo tiene 4 series mensuales en 2025–2026; el norte tiene además la
ocupación semanal.

\begin{center}
\small
\begin{tabularx}{\linewidth}{>{\raggedright\arraybackslash}p{0.3\linewidth}Y}
\encabezado \h{Opción} & \h{Por qué sí o por qué no} \\
\textbf{Pesos iguales con sensibilidad de ±50\,\% (elegida)} & Fácil de defender; no premia al norte por tener más series \\
\filagris Componentes principales & Pesos ``objetivos'', pero difíciles de explicar e inestables con 19 meses del sur \\
Criterio del equipo & Subjetivo; cada número habría que justificarlo \\
\bottomrule
\end{tabularx}
\normalsize
\end{center}
\textbf{Resultado}: moviendo cada peso a 0.5 y 1.5, cambia de estado entre el 0.6\,\% y el 6.0\,\% de los lugar-mes.
\end{opciones}

\begin{opciones}[4.2 · Los cortes entre estados · 28-sep-2026]
\textbf{Elegida}: \textbf{percentiles comunes} a todo el estado: debajo de la mediana, tranquilo; arriba del percentil 90,
saturado.

\textbf{Descartados}: terciles de cada lugar (todo lugar saldría ``saturado'' un tercio del tiempo aunque nunca se llene:
Cozumel promedia 53\,\%) y grupos k-means (los grupos pueden no significar nada y cambian al agregar datos).
\end{opciones}

\begin{opciones}[4.3 · Predecir con aprendizaje de máquina · 28-sep-2026]
\textbf{Lo que dijo Brandon}: \emph{``con los datos que tenemos inferimos los futuros haciendo ML''}.

\textbf{Consecuencia}: un clasificador del estado del mes siguiente (logística, Random Forest y Gradient Boosting) y una
cadena de Markov. El Radar predice \emph{estados}; los \emph{niveles} son de la Fase 5.
\end{opciones}

\begin{opciones}[4.4 · La escala de las variables · 29-sep-2026]
\textbf{Evidencia presentada}: en mín–máx, el mes típico de visitantes INAH por mil habitantes cae en 0.02, el del Tren Maya
en 0.06, el de cruceros en 0.11, el de ocupación en 0.70 y el de Belice en 0.74.

\textbf{Elegida}: \textbf{mín–máx} con mínimo y máximo comunes. \textbf{Descartados}: percentil común (pierde las
proporciones) y mín–máx sobre logaritmo (más difícil de explicar).
\end{opciones}

\begin{opciones}[4.5 · La prueba de validez que falló: opción D · 29-sep-2026]
\textbf{Lo que pasó}: con solo SITUR-Q, Cancún salía \textbf{``tranquilo''} en julio de 2026 (0.025), porque SITUR-Q dejó de
publicar su ocupación; DataTur marcaba 68.8\,\%. Además, Chetumal (0.57) salía más presionado que Cancún (0.27) en 2024 por
los cruces de Belice.

\begin{center}
\small
\begin{tabularx}{\linewidth}{>{\raggedright\arraybackslash}p{0.28\linewidth}rrrrY}
\encabezado \h{Variante} & \h{Cancún} & \h{Isla M.} & \h{Chetumal} & \h{Ruta} & \h{Problema} \\
A. Mín–máx, solo SITUR-Q & 0.02 & — & 0.39 & 0.14 & El norte ``desaparece'' \\
B. + ocupación DataTur & 0.28 & 0.45 & 0.39 & 0.14 & Chetumal arriba de Cancún \\
C. Percentil + DataTur & 0.44 & 0.25 & 0.48 & 0.71 & La Ruta arriba de Cancún \\
\textbf{D. B sin componentes de un solo lugar} & \textbf{0.28} & \textbf{0.45} & \textbf{0.05} & \textbf{0.14} & Quiebre de
2025 \\
\bottomrule
\end{tabularx}
\normalsize
\end{center}

\textbf{Elegida por Brandon: D}. La ocupación de los 4 lugares que mide DataTur viene de DataTur en toda su historia, y un
componente entra solo si lo tienen al menos 2 lugares (sale Belice, que solo tiene Chetumal).
\end{opciones}

\begin{opciones}[4.6 · Qué índice se usa para predecir y cuál se publica · 29-sep-2026]
\textbf{El problema}: en 2025 Chetumal ``cayó'' de 0.54 a 0.05 porque perdió la ocupación, no porque llegara menos gente. Un
modelo aprendería ese cambio de datos como si fuera presión.

\textbf{Elegida}: el \textbf{índice comparable}: cada lugar usa solo las medidas que tiene hoy, y solo en los meses en que
las tiene todas. Se usa para predecir \emph{y} para publicar (un solo índice: con dos, Cancún salía ``tranquilo'' en uno y
``concurrido'' en el otro el mismo mes).

\textbf{Descartados}: entrenar solo con 2025–2026 (19 meses por lugar) y entrenar todo con una bandera ``tiene ocupación''
(el modelo aprendería el cambio de datos).
\end{opciones}

\begin{opciones}[4.7 · Cómo elegir el modelo · 29-sep-2026]
\textbf{Criterio de Brandon}: \textbf{el que más acierta y, a igualdad, el que más cambios anticipa}, porque la regla
anti-colapso necesita anticipar.

\textbf{Descartado}: elegir por F1-macro, la métrica del plan.

\textbf{Resultado vigente} (tras la auditoría): la regresión logística, con 130 de 156 aciertos y 8 de 30 cambios
anticipados, contra 130 y 6 de Random Forest y 126 y 0 de ``igual que este mes''. Antes de la auditoría, el mismo criterio
daba Random Forest: \textbf{se respeta el criterio, no el nombre del modelo}, y el código elige solo.

\textbf{Lectura honesta}: el estado se repite el 81\,\% de los meses, así que el modelo agrega poco; su valor está en los
cambios que anticipa. En los 5 lugares hubo solo 2 cambios en 48 casos y no anticipó ninguno: no se puede afirmar que
anticipe saturaciones del sur.
\end{opciones}

\begin{opciones}[4.8 · Markov: ¿semanal o mensual? · 29-sep-2026]
\textbf{Elegida}: \textbf{semanal, solo el norte} (7 centros de DataTur, 239 semanas, 1,666 transiciones). Es lo que marca el
plan (Markov en semanas, el Pronóstico en meses) y sirve a la campaña: cuando el norte se va a saturar, es el momento de
ofrecer el sur.

\textbf{Descartados}: mensual para todos (repetía al clasificador y se encimaba con la Fase 5) y las dos.

\textbf{Resultado}: mejor puntaje de Brier que la persistencia a 1, 4 y 8 semanas (0.200 contra 0.225; 0.389 contra 0.484;
0.490 contra 0.676).
\end{opciones}

\section*{La auditoría del 29 de septiembre}
Brandon pidió: \emph{``vuelve a revisar la fase 4 auditando todo 1 al 4 y si todo está bien según el plan sigamos''}. Se
revisó punto por punto contra el plan, se volvió a correr todo desde cero (las 8 tablas del Radar salieron idénticas) y una
prueba nueva comprobó 20 cifras de los documentos contra el código. Faltaban tres puntos del plan:

\begin{opciones}[4.9 · Lo que faltaba y lo que decidió Brandon · 29-sep-2026]
\begin{center}
\small
\begin{tabularx}{\linewidth}{>{\raggedright\arraybackslash}p{0.28\linewidth}Y}
\encabezado \h{Hallazgo} & \h{Decisión y resultado} \\
El índice no tenía ``llegadas por cuarto'' ni ``densidad de oferta'' & Se agregan llegadas por cuarto (tren + cruceros). La
densidad no entra: es un solo corte en el tiempo y no señala meses. Se le mostró a Brandon que con cruceros el componente
queda dominado por los puertos (Mahahual: 432 por cuarto al mes) y decidió mantener tren + cruceros, declarándolo \\
\filagris Faltaba el agrupamiento de los centros del país & Ward con 55 centros completos; $k=2$ por silueta (0.450).
Descartado $k=7$ (0.317) \\
DataTur mide menos que SITUR-Q & Se declara, sin ajustar: Cozumel $-14.5$, Isla Mujeres $-19.3$, Cancún $-1.6$ puntos. El
error es conservador (el norte se ve menos lleno). Descartado calibrar, que volvería estimación la ocupación del norte \\
\filagris El sesgo del modelo & Medido y declarado (decisión 4.7) \\
\bottomrule
\end{tabularx}
\normalsize
\end{center}
\textbf{Cifras vigentes}: cortes del índice comparable 0.186 y 0.756; en julio de 2026 los 5 lugares están tranquilos y
Cancún (0.193) concurrido.
\end{opciones}

\begin{error}
\begin{itemize}
\item \textbf{La ``ocupación'' de 2,560,068}: SITUR-Q trae tres renglones por mes (cuartos disponibles, ocupados y el
  porcentaje) y el panel los sumaba. Ahora se recalcula ocupados entre disponibles (Chetumal 2024: 58.0\,\%) y el panel se
  detiene si una variable trae dos renglones el mismo mes.
\item \textbf{Visitantes contados dos veces}: la serie ``zonas arqueológicas'' de SITUR-Q es el INAH sumado por destino
  (Chetumal 2025: $39{,}459=11{,}017+21{,}850+6{,}592$). Se usa solo el INAH, zona por zona.
\end{itemize}
\end{error}
